\documentclass[10pt,a4paper]{article}

\usepackage[T1]{fontenc}
\usepackage[utf8]{inputenc}
\usepackage{times}
\usepackage{microtype}
\usepackage[margin=0.85in]{geometry}
\usepackage{cite}
\usepackage{url}

\setlength{\parindent}{0pt}
\setlength{\parskip}{4pt}

\title{Literature Review: Multi-Page Visually Rich Document Understanding with Agentic and Knowledge-Grounded Retrieval}
\author{Lewei Xu}
\date{May 2026}

\begin{document}
\maketitle

\section{Introduction}

Visually rich documents are a primary medium for institutional, scientific, financial, and legal information. Unlike plain text corpora, these documents encode meaning through text, layout, tables, figures, charts, page order, section hierarchy, and cross-references. A question about a financial report may require reading a table on one page, locating a definition in a footnote on another, and applying domain knowledge about reporting conventions. A question about a scientific article may require matching a result statement with a figure panel, caption, method section, and appendix table. These cases motivate visually rich document understanding (VRDU), where the task is not only to read words, but to ground information in document structure and visual context.

Early document visual question answering systems largely focused on single-page settings. This restriction made the task tractable while preserving important difficulties such as optical character recognition (OCR), spatial layout modelling, and multimodal fusion. Recent MLLM-based VRDU work has substantially improved single-page document understanding through stronger visual encoders, layout-aware pretraining, OCR-free reading objectives, and instruction tuning \cite{huang2022layoutlmv3,hu2024mplug,ding2025survey}. However, many real documents exceed a single page, and their meaning is often distributed across non-adjacent locations. Multi-page visually rich document understanding (MP-VRDU) therefore introduces additional problems of document-scale evidence selection, cross-page reasoning, and long-context control.

The central difficulty in MP-VRDU is that document length changes the nature of the problem. A single-page model can often process all visual and textual evidence jointly. A multi-page model must decide what to retain, retrieve, compress, revisit, or ignore. Directly encoding every page produces rapidly growing OCR token sequences and visual token sequences. Retrieval reduces the context burden, but missed pages can become unrecoverable. Agentic systems improve adaptability by iteratively searching and inspecting evidence, but they introduce tool latency, dependence on external parsers, and less reproducible decision processes. These trade-offs make MP-VRDU a useful test case for broader questions in grounded multimodal reasoning.

This review synthesises the literature that motivates research on multi-agent decomposition with retrieval-augmented grounding for unfamiliar document domains. The aim is not to enumerate every model, but to explain the progression from single-page document understanding to document-scale retrieval and agentic reasoning. The review first outlines the foundations of MP-VRDU, then organises recent methods into architectural trends, evidence modelling strategies, retrieval-based systems, and agentic pipelines. It then identifies a gap that is only partially addressed by current work: most systems ground answers in the target document and in benchmark-specific training distributions, while unfamiliar domains may require external domain knowledge such as terminology, schemas, regulatory rules, or scientific background. This motivates research on systems that can combine document-grounded evidence with domain-knowledge retrieval while preserving traceability and verification.

\section{From Single-Page DocVQA to MP-VRDU}

Document VQA emerged from the observation that document images require models to read, localise, and reason over text in visual context. This differs from natural-image VQA, where scene objects dominate, and from text-only QA, where spatial arrangement is usually absent. Document pages contain dense text, repeated fields, tables, captions, headers, signatures, and visual separators. Layout-aware models such as LayoutLMv3 encode OCR tokens, bounding boxes, and image patches within a unified pretraining framework, thereby improving the alignment between textual and spatial evidence \cite{huang2022layoutlmv3}. OCR-free systems such as Donut and later document-focused MLLMs instead process document images directly and learn text recognition or structure recovery as part of visual-language modelling \cite{kim2021donut,hu2024mplug}. These advances made single-page VRDU increasingly effective, but they did not remove the multi-page bottleneck.

The MP-DocVQA setting formalised the requirement to answer questions over documents rather than isolated pages. Hi-VT5 introduced a hierarchical multimodal transformer that encodes each page and aggregates page-level information for answer generation \cite{tito2023hierarchical}. This architecture reflects a key insight: page boundaries are not incidental formatting details, but part of the document structure. Preserving page identity lets the model distinguish repeated fields, locate evidence, and reason about answer provenance. Yet page-level aggregation also compresses information before full document reasoning, which can remove fine-grained evidence needed for later comparisons.

Subsequent hierarchical systems focused on stronger document-level exchange. GRAM adds document-global tokens that mediate communication across page representations, allowing information to flow beyond isolated page encoders \cite{blau2024gram}. Recurrent-memory approaches pass memory tokens from page to page, reducing the quadratic cost of full concatenation while preserving some cross-page state \cite{dong2024multi}. Arctic-TILT extends OCR-based document modelling to much longer inputs through chunked sparse attention and layer-wise multimodal fusion, showing that careful attention design and engineering can substantially increase the number of processable pages \cite{borchmann2025arctic}. These methods demonstrate that long documents can be modelled more directly than early single-page systems allowed, but they also reveal persistent limits. Summaries, global tokens, sparse attention masks, and recurrent memories all decide what information survives. Fine-grained evidence may be lost if it is compressed before the model knows its relevance to the question.

The transition from single-page to multi-page VRDU also exposes long-context weaknesses known from language modelling. Long inputs can dilute attention, and models often underuse evidence placed in the middle of a context window \cite{liu2024lost}. The visual setting amplifies this issue. Each page may require hundreds or thousands of visual tokens at useful resolution, and text-heavy pages may require OCR tokens, coordinates, and crop-level information. A model that consumes all pages jointly therefore faces a tension between coverage and fidelity. Low-resolution images preserve global layout but lose small text. High-resolution images preserve details but exceed context budgets. OCR gives compact text but can propagate recognition and ordering errors. MP-VRDU methods are best understood as different responses to this tension.

Benchmarks have also broadened the evaluation target. DUDE evaluates document understanding across diverse document types and emphasises the limitations of narrow benchmark assumptions \cite{van2023document}. MMLongBench-Doc and LongDocURL extend evaluation toward long-context document understanding with visualisations, locating requirements, and reasoning over many pages \cite{ma2024mmlongbench,deng2025longdocurl}. M3DocVQA further stresses multi-page and multi-document understanding, requiring systems to retrieve and reason across a wider evidence space \cite{cho2025m3docvqa}. These datasets show that MP-VRDU is not just a scaled-up DocVQA task. It requires document navigation, evidence attribution, robustness to heterogeneous layouts, and reasoning over information that may be separated by page boundaries, modality boundaries, or domain conventions.

\section{Architectural Trends in MP-VRDU}

Recent MP-VRDU methods can be organised into four broad architectural families: page-to-document transformers, backbone-centric MLLM adaptations, retriever-generator systems, and agentic pipelines. These families are not mutually exclusive, but they clarify the main design decisions. The first two families try to make the document fit inside a model. The latter two make the model decide which parts of the document to inspect.

Page-to-document transformers extend document encoders by preserving page-local representations and adding document-level aggregation. Hi-VT5 uses page-wise encoding followed by answer generation over aggregated evidence \cite{tito2023hierarchical}. GRAM strengthens cross-page interaction through global tokens and local-global attention \cite{blau2024gram}. RM-T5 uses recurrent memory to carry information across pages \cite{dong2024multi}. Arctic-TILT uses sparse long-context attention over OCR streams and repeated text-image-layout fusion \cite{borchmann2025arctic}. The strength of this family is its explicit inductive bias for page structure. It can model all or most pages in one document-level pass, which supports comparison and aggregation when evidence is distributed. The limitation is that computation generally grows with document length, and the information exchange mechanism may become the bottleneck. If relevant evidence is compressed away or never attended to, later reasoning cannot recover it.

Backbone-centric MLLM adaptations instead treat MP-VRDU as a problem of adapting a general multimodal model to long, text-rich documents. Some approaches introduce layout-aware tokens, document-specific projectors, or instruction tuning to align page images, OCR, and layout with an LLM backbone \cite{zhu2025simple,tanaka2024instructdoc}. Others focus on reducing visual tokens so that more pages can remain in context. mPLUG-DocOwl2 compresses high-resolution document images into a smaller number of visual tokens while preserving layout-aware information \cite{hu2025mplug}. DocSLM combines visual, OCR, and layout cues into fixed-size page representations and uses streaming abstention to process long documents under limited hardware \cite{hannan2025docslm}. Docopilot trains native document-level MLLMs on large multi-page document corpora and engineering optimisations for long-context processing \cite{duan2025docopilot}. These systems benefit from the instruction-following and multimodal reasoning capacity of foundation models, but they still require document-scale evidence to survive inside the backbone context. Their failure modes therefore centre on visual token growth, compression loss, context limits, and poor recovery when the first pass overlooks evidence.

Retriever-generator systems decompose MP-VRDU into evidence selection and answer generation. The retriever identifies pages, patches, chunks, figures, or graph nodes, and the generator answers from this reduced context. RAG-DocVQA showed that retrieval can improve multi-page document VQA by avoiding full-document concatenation and by passing only selected textual or visual evidence to a downstream model \cite{lopez2025enhancing}. M3DocRAG demonstrated the importance of multimodal retrieval for multi-page and multi-document settings, using visual document retrievers such as ColPali or ColQwen to keep visual and textual evidence coupled at page level \cite{cho2024m3docrag,faysse2024colpali}. VDocRAG adapts LVLMs toward document retrieval through representation compression objectives, reducing the gap between generative visual-language pretraining and retrieval needs \cite{tanaka2025vdocrag}. AVIR and related systems address the rigidity of fixed top-$k$ retrieval by adaptively selecting pages according to score distributions \cite{li2025avir}. Retriever-generator methods scale well, provide partial evidence localisation, and can plug into strong external MLLMs. Their core weakness is retrieval recall. If the retriever misses a necessary page, crop, or table, the generator is usually forced to infer without evidence.

Agentic pipelines extend retrieval by making evidence acquisition iterative. Rather than retrieving once, an agent decomposes the question, searches, inspects, revises, and verifies. MDocAgent uses multiple agents over textual and visual retrieval branches, separating preliminary answering, critical cue extraction, modality-specific analysis, and final synthesis \cite{han2025mdocagent}. DocAgent constructs a document outline and lets an actor use search, section, image, page, and table tools, with a reviewer checking the answer through complementary evidence \cite{sun2025docagent}. DocLens combines page navigation, element localisation, answer sampling, and adjudication to improve evidence-page recall and reduce hallucination on long documents \cite{zhu2025doclens}. Doc-React formalises iterative retrieval as a ReAct-style process that decomposes residual information needs and retrieves new multimodal evidence \cite{wu2025doc}. ViDoRAG adds hybrid retrieval, adaptive evidence selection, and multi-scale agent inspection \cite{wang2025vidorag}. Doc-V* trains an OCR-free visual agent to inspect document thumbnails, fetch high-resolution pages, maintain working memory, and answer after sequential evidence aggregation \cite{zheng2026doc}. These systems point toward a shift from document encoding to document navigation.

The architectural trend is therefore clear. Early approaches ask how to represent many pages inside a model. Later approaches ask how to select and inspect the right evidence. Recent agentic methods ask how a system should decide what to inspect next. This progression is especially relevant for unfamiliar domains. A model may not know which fields, tables, abbreviations, or external definitions are important until it reasons about the question. Fixed retrieval and single-pass encoding are therefore poorly matched to open-ended document analysis, where the evidence need may evolve during reasoning.

\section{Modelling Evidence Across Text, Vision, Layout, and Pages}

MP-VRDU systems differ not only in architecture, but also in how they represent evidence. Text is the most compact and searchable modality. OCR text, embedded PDF text, generated summaries, and synthetic query nodes all provide efficient retrieval handles. Hierarchical transformers use OCR tokens and bounding boxes as structured page inputs \cite{tito2023hierarchical,borchmann2025arctic}. RAG systems chunk OCR or parsed text, embed it, and retrieve semantically relevant passages \cite{lopez2025enhancing,shin2025multidocfusion}. Agentic systems use text as a navigation layer, searching outlines, section titles, captions, and extracted snippets before deciding whether visual inspection is necessary \cite{sun2025docagent,zhu2025doclens}. The limitation is that text alone is not a faithful document representation. OCR may misread symbols, reorder multi-column layouts, lose table structure, and omit chart semantics. Extracted text also strips visual context, which can be decisive in forms, slides, figures, diagrams, and scanned reports.

Visual evidence preserves what text extraction may lose. Purely visual retrievers represent pages as image embeddings, allowing retrieval over screenshots or rendered pages without relying on OCR at inference \cite{cho2024m3docrag,tanaka2025vdocrag}. High-resolution document MLLMs preserve text, layout, and figures directly from pixels, but must compress visual tokens aggressively to support multiple pages \cite{hu2025mplug,hannan2025docslm}. Agentic systems often separate cheap visual overview from expensive visual reading. Thumbnail grids or page previews can guide coarse navigation, while high-resolution crops or page fetches support final evidence extraction \cite{zheng2026doc,wang2025vidorag}. This mirrors human document use: readers skim the document structure before zooming into pages, tables, or figures.

Layout connects text and vision. At page level, layout includes coordinates, reading order, regions, boxes, and relative positions. At document level, it includes headings, section hierarchy, page order, cross-references, tables continuing across pages, figure-caption relations, and appendix links. The multi-page setting makes document-level layout especially important. MultiDocFusion reconstructs hierarchical document trees before chunking, showing that structure-aware segmentation can preserve context that fixed-size chunks break \cite{shin2025multidocfusion}. MHier-RAG builds multi-granularity indexes over flattened page chunks and cross-page summaries, allowing retrieval to operate at page and document levels \cite{gong2025mhier}. MLDocRAG constructs a multimodal chunk-query graph by generating query-answer anchors for different element types, turning heterogeneous document units into searchable graph nodes \cite{zhang2026mldocrag}. MoLoRAG builds a multimodal page graph and traverses it using both semantic similarity and logic-aware relevance \cite{wu2025molorag}. These methods show that layout is not only geometry. It is also a retrieval structure.

Fusion strategies determine when modalities interact. Early fusion combines OCR, visual patches, layout, and question tokens before or during encoding. This supports fine-grained alignment but scales poorly across many pages. Late fusion keeps modalities separate during retrieval and combines them only after selecting candidate evidence. This is more scalable, but it can miss cross-modal dependencies if the retrieval stage is weak. Staged fusion offers a middle path: fuse locally within pages or regions, compress or index the result, then perform selective document-level reasoning. Many current systems approximate this pattern through page summaries, visual embeddings, working memory, graph nodes, or retrieved multimodal chunks \cite{hannan2025docslm,cho2024m3docrag,gong2025mhier}. For unfamiliar domains, staged fusion is attractive since it preserves page-local grounding while allowing retrieval and agents to adaptively choose what evidence needs deeper inspection.

Evidence modelling therefore shapes the ceiling of downstream reasoning. A system cannot ground an answer in a table if the table has been flattened into unordered text. It cannot reason over a figure if the visual retriever only stores text summaries. It cannot compare cross-page facts if page identity is lost. The most robust MP-VRDU systems are likely to maintain multiple evidence views: searchable text, visual page representations, layout-aware structure, and traceable links back to source regions.

\section{RAG and Agentic Reasoning for Long Document Understanding}

Retrieval-augmented generation has become central to MP-VRDU as a response to context limits. In its simplest form, RAG retrieves relevant document text or pages and gives them to a generator. This reduces token load and provides a natural evidence trace. However, document RAG is harder than text RAG. Relevant evidence may be visual rather than textual, may be located through layout rather than semantic similarity, or may require domain knowledge to recognise. A keyword query may retrieve a page containing the requested term but miss a table that defines the value. A visual retriever may find similar-looking pages but miss logical relevance. A text embedder may fail when field names are unfamiliar, abbreviated, or domain-specific.

MP-VRDU RAG systems respond to these problems through richer retrieval units and stronger retrieval objectives. RAG-DocVQA evaluates textual and visual RAG variants, showing that retrieval can outperform concatenate-all-pages baselines when the retriever is well aligned with the task \cite{lopez2025enhancing}. M3DocRAG uses multimodal page retrieval to preserve visual evidence and avoid reliance on OCR-only text \cite{cho2024m3docrag}. VDocRAG trains document visual representations specifically for retrieval, rather than treating generative LVLM embeddings as sufficient \cite{tanaka2025vdocrag}. AVIR adapts the number of selected pages to the query and score distribution, reflecting that different questions need different evidence budgets \cite{li2025avir}. These works collectively show that retrieval is not a preprocessing convenience. It is a central modelling choice that determines what the generator can know.

Graph and structure-aware RAG further address the limits of flat retrieval. Knowledge graph prompting connects passages or document units through graph relations, supporting multi-hop reasoning when direct semantic similarity is insufficient \cite{wang2024knowledge}. MHier-RAG retrieves at multiple granularities so that a hit on a small text block can return its parent page and surrounding visual context \cite{gong2025mhier}. MLDocRAG links generated query nodes to multimodal chunks, allowing the system to retrieve likely answer-bearing chunks through synthetic semantic anchors \cite{zhang2026mldocrag}. MoLoRAG uses graph traversal to combine semantic relevance with logical relevance over multimodal page nodes \cite{wu2025molorag}. These systems are important for unfamiliar domains since domain-specific terms and implicit relationships may not be captured by surface similarity alone.

Agentic reasoning adds a control layer above retrieval. A fixed RAG pipeline assumes that the initial query is sufficient and that one retrieval pass can identify all evidence. In MP-VRDU, both assumptions are weak. Questions may require decomposition, such as finding a definition, then locating a value, then comparing it with a related figure. Evidence may be scattered across text and visual elements. The first retrieved page may reveal new terms that should guide the next search. Agentic pipelines therefore let an LLM or VLM decide the next action based on observations. This action may be searching text, retrieving images, fetching a page, cropping a table, asking a specialist model to read a figure, or verifying a candidate answer.

Multi-agent decomposition is one way to make this control process manageable. MDocAgent separates retrieval, critical cue extraction, modality-specific analysis, and synthesis, reducing the burden on a single model \cite{han2025mdocagent}. M2RAG combines text retrieval, visual retrieval, visual extraction, and modal fusion agents, treating multimodal RAG as a collaboration among specialised components \cite{duan2025m2rag}. DocLens decomposes long-document QA into page navigation, element localisation, answer sampling, and adjudication \cite{zhu2025doclens}. ViDoRAG uses coarse-to-fine multi-scale inspection, where agents progressively narrow a visual evidence space \cite{wang2025vidorag}. These systems suggest that agent roles are useful when the task itself has separable subproblems: locating evidence, reading visual content, interpreting structure, judging sufficiency, and producing an answer.

Single-agent iterative systems provide a contrasting design. Doc-React uses a ReAct-style loop that reformulates sub-queries and retrieves additional evidence according to remaining uncertainty \cite{wu2025doc}. SimpleDoc uses dual-cue retrieval and iterative refinement with a working memory, showing that a well-designed single agent can compete with heavier multi-agent pipelines \cite{jain2025simpledoc}. Doc-V* trains an OCR-free visual agent to inspect thumbnails and fetch pages, combining learned action selection with visual working memory \cite{zheng2026doc}. DocDancer trains an open-source document QA agent from synthetic tool-use trajectories, reducing reliance on prompt-only closed-source orchestration \cite{zhang2026docdancer}. These systems show that the key issue is not the number of agents by itself, but whether the system can adapt its evidence acquisition process.

The main limitation of current agentic MP-VRDU systems is grounding discipline. Agents can call tools, retrieve evidence, and produce intermediate reasoning, but their final answers may still blend document evidence, model priors, and unsupported assumptions. Multi-agent systems can also amplify inconsistency if agents disagree without reliable adjudication. Tool outputs may be noisy, and expensive visual inspections may be used inefficiently. These concerns become sharper when external domain knowledge is added. The system must then distinguish what the document states from what external sources imply, and it must decide when external knowledge is allowed to clarify, constrain, or challenge document-grounded evidence.

\section{Domain Knowledge Grounding for Unfamiliar Documents}

Most MP-VRDU benchmarks evaluate whether a model can find and read evidence inside the target document. This is necessary, but not sufficient for unfamiliar document domains. Real documents often rely on conventions that are external to the page. Financial reports use accounting terms and reporting schemas. Medical forms use abbreviations and clinical units. Legal contracts depend on defined terms and statutory concepts. Scientific articles assume background knowledge about methods, datasets, metrics, and notation. A reader may need this knowledge to recognise which evidence matters, interpret table headings, disambiguate acronyms, or verify whether an answer is plausible.

Current MP-VRDU systems address domain unfamiliarity mainly through larger backbones, broader instruction tuning, or better document retrieval. Backbone-centric MLLMs inherit general knowledge from pretraining, but this knowledge is implicit and difficult to cite. Retriever-generator systems ground answer generation in document evidence, but they usually retrieve only from the target document or corpus. Agentic systems can use tools, but most document agents focus on parsers, search, page images, and tables rather than external knowledge sources \cite{sun2025docagent,zhu2025doclens,zhang2026docdancer}. This leaves a gap between document-grounded QA and domain-grounded interpretation.

External domain-knowledge RAG can fill part of this gap. Instead of retrieving only document pages, the system can retrieve definitions, schema descriptions, regulations, ontology entries, technical background, or examples from a curated knowledge source. This retrieval should not replace document evidence. It should ground interpretation of the document. For example, external knowledge may explain that a field name is an abbreviation, that a regulatory threshold determines which table row is relevant, or that a metric in a scientific paper has a standard definition. The answer should still cite the document page or region that contains the case-specific value.

The challenge is controlling the relationship between internal and external evidence. A document-grounded answer can be wrong if the system lacks domain knowledge. An externally grounded answer can also be wrong if the retrieved knowledge is irrelevant, outdated, or conflicts with the document. Domain knowledge should therefore be treated as contextual support for interpretation and verification, not as an unconstrained source of answers. This requires evidence typing. The system should know whether a piece of retrieved information is a document fact, a domain definition, a schema rule, or a reasoning intermediate. It should also preserve source links so that users can distinguish document-grounded claims from background assumptions.

Graph-based and agentic methods provide useful building blocks for this direction. Knowledge graph prompting shows that graph structures can improve multi-hop question answering by making relations explicit \cite{wang2024knowledge}. MoLoRAG shows that graph traversal can combine semantic and logical relevance over document pages \cite{wu2025molorag}. MHier-RAG and MLDocRAG show how document chunks, summaries, and generated queries can be organised into retrieval structures beyond flat vectors \cite{gong2025mhier,zhang2026mldocrag}. Agentic systems show how retrieval can be guided by intermediate reasoning and tool feedback \cite{wu2025doc,wang2025vidorag,zheng2026doc}. A natural next step is to connect these ideas: document graph nodes, visual evidence, and external domain-knowledge nodes can be retrieved and inspected by agents that maintain separate evidence channels.

This direction remains underdeveloped in MP-VRDU. Existing systems generally optimise for answer accuracy on known benchmarks, not for robust transfer to unfamiliar document domains. They also rarely evaluate whether external knowledge improves interpretation while preserving document faithfulness. A system for unfamiliar domains should therefore be judged on more than answer accuracy. It should retrieve relevant document evidence, retrieve relevant external domain knowledge when needed, avoid using external knowledge when the document is sufficient, resolve conflicts conservatively, and expose a trace that distinguishes source types.

\section{Datasets, Evaluation, and Remaining Gaps}

The choice of benchmark shapes the perceived progress of MP-VRDU. MP-DocVQA established the core setting of answering questions over multi-page documents, but many questions remain extractive and depend on locating answer-bearing pages \cite{tito2023hierarchical}. DUDE expands the range of document types and highlights robustness across more realistic forms of document understanding \cite{van2023document}. MMLongBench-Doc focuses on long-context multimodal document understanding with visualisations, stressing models that must process many pages and figures \cite{ma2024mmlongbench}. LongDocURL adds locating and reasoning requirements over long multimodal documents \cite{deng2025longdocurl}. M3DocVQA extends evaluation to multi-modal, multi-page, and multi-document settings \cite{cho2025m3docvqa}. These benchmarks are important, but they still only partially capture the demands of unfamiliar-domain document reasoning.

A first gap is evidence sufficiency. Many systems report answer metrics such as ANLS or accuracy, but a correct answer may be produced without robust evidence grounding. Conversely, a system may retrieve the correct page but fail to extract the answer due to OCR or visual reading errors. Evidence-page recall, region localisation, and citation fidelity are therefore important for interpreting model performance. DocR1 explicitly trains evidence-page prediction with reinforcement learning, indicating growing recognition that answer generation and evidence localisation should be coupled \cite{xiong2026docr1}. Agentic systems such as DocLens also prioritise evidence-page recall before answer synthesis \cite{zhu2025doclens}. Future evaluation should extend this idea to document regions, table cells, figures, and external knowledge sources.

A second gap is domain transfer. Many MP-VRDU systems are trained or tuned on benchmark-like data. Retrieval models may require task-specific fine-tuning, and agent trajectories may be generated from known document distributions. This limits confidence in unfamiliar domains where vocabulary, layout conventions, and evidence forms differ. Training-free or training-light systems are attractive for this reason, but prompt-only agentic systems can be brittle and expensive. A useful research direction is therefore not simply to avoid training, but to design modular systems whose retrieval, parsing, visual inspection, and domain-knowledge grounding can transfer with minimal new labelled data.

A third gap is adaptive reasoning under retrieval failure. Fixed RAG systems are efficient, but they provide little recourse when the initial evidence set is incomplete. Agentic systems address this through iterative search, but current designs often depend on hand-written prompts, closed-source models, or learned trajectories tied to specific tool sets. More general systems need explicit policies for when to reformulate a query, switch modality, retrieve external knowledge, inspect a crop, ask a verifier, or stop. These decisions should be evaluated not only by final answer accuracy, but by cost, latency, evidence trace quality, and robustness to misleading retrieval results.

A fourth gap is conflict management. MP-VRDU systems often treat retrieved context as uniformly reliable. In unfamiliar domains, retrieved document evidence, parser output, visual readings, and external knowledge may disagree. OCR may misread a value; a table parser may misalign columns; an external definition may not apply to the document's local terminology. Multi-agent decomposition can help by assigning specialised roles for retrieval, reading, domain interpretation, and verification, but role specialisation alone does not guarantee faithful resolution. Systems need mechanisms to classify evidence sources, preserve provenance, and prefer document-grounded claims when external knowledge is only explanatory.

These gaps motivate a research direction centred on multi-agent, retrieval-augmented MP-VRDU for unfamiliar domains. The target is not a monolithic model that encodes every page, nor a fixed RAG pipeline that retrieves once and answers. A more suitable framework would decompose the task into evidence search, visual and textual reading, layout-aware structure use, external domain-knowledge retrieval, and verification. It would maintain separate but linked evidence channels for document facts and background knowledge. It would use agents or agent-like controllers to decide when the current evidence is insufficient and which tool or retrieval source should be used next. This direction follows naturally from the literature: MP-VRDU has moved from full-context modelling toward selective navigation, and unfamiliar-domain use cases require that navigation to extend beyond the document while remaining grounded in it.

\section{Conclusion}

The literature on MP-VRDU shows a clear progression. Single-page document understanding established methods for reading text, layout, and visual evidence. Multi-page benchmarks exposed the limits of page-local modelling and full-document concatenation. Hierarchical transformers, long-context MLLMs, retriever-generator systems, and agentic pipelines each address part of the document-scale problem. Hierarchical and backbone-centric methods preserve joint document context but strain memory and compression. Retriever-generator methods scale more effectively but depend on evidence recall. Agentic systems adaptively seek evidence, but introduce higher cost, tool dependence, and new grounding risks.

The remaining challenge is not merely to answer questions over longer documents. It is to answer them faithfully when evidence is distributed, multimodal, and embedded in unfamiliar domain conventions. Current methods mostly ground answers in the target document and in knowledge implicit in model parameters. A more robust direction is to combine document-grounded retrieval with external domain-knowledge retrieval, controlled by agents that can decompose questions, inspect evidence, verify claims, and preserve provenance. Such systems would better support users who need to understand unfamiliar visually rich documents, while giving future researchers a clearer basis for comparing document navigation, evidence grounding, and domain transfer.

\bibliographystyle{ieeetr}
\bibliography{custom}

\end{document}
